% ============================================================
% ch20.tex —— 第 20 章 给素数画人口普查图
% ============================================================
\chapter{给素数画人口普查图}

第 4 章挂起的巨额欠条，现在开始偿还。问题原样回收：素数有无穷多个（欧几里得），但它们分布的\textbf{节奏}是什么？这一章不证明任何东西，只做人口学家的事：\textbf{给素数画普查图}——而图上会浮现出一个谁都想不到的公式，它里面居然藏着 $\ln$、藏着 $\mathrm{e}$。素数是整数的事，对数是连续增长的事，它们怎么会认识？认识的过程，本身就是"解析数论"这门学科的诞生记。

\section{亲手数素数}

先自己动手，数出第一批数据（比读一百页教材都管用）。工具是第 4 章的埃拉托斯特尼筛法：写下 $1$ 到 $100$，划掉 $1$（不是素数）；留下 $2$、划掉 $2$ 的倍数；留下 $3$、划掉 $3$ 的倍数（$4, 6, 8\dots$ 已划）；留下 $5$、划掉 $5$ 的倍数；留下 $7$、划掉 $7$ 的倍数——到此为止（$7^2 = 49 < 100 < 121 = 11^2$，试到 $\sqrt{100}$ 即可，第 3 章的法则）。数一数幸存者：$25$ 个。

继续数（这道苦工外包给计算机，但值得亲手数到 $200$ 找感觉）：$200$ 以内 $46$ 个。批量数据入表：

\begin{center}
\begin{tabular}{ccc}
\toprule
范围 $x$ & 素数个数 $\pi(x)$ & 占比 $\pi(x)/x$ \\
\midrule
$10$ & $4$ & $40\%$ \\
$100$ & $25$ & $25\%$ \\
$1\,000$ & $168$ & $16.8\%$ \\
$10\,000$ & $1\,229$ & $12.3\%$ \\
$100\,000$ & $9\,592$ & $9.6\%$ \\
$1\,000\,000$ & $78\,498$ & $7.85\%$ \\
$10^9$ & $50\,847\,534$ & $5.08\%$ \\
$10^{18}$ & $24\,739\,954\,287\,740\,860$ & $2.47\%$ \\
\bottomrule
\end{tabular}
\end{center}

记号：$\pi(x)$（素数计数函数；与圆周率 $\pi$ 无关，只是重名——历史的选择，忍一忍）。三句读后感：
\begin{enumerate}
\item \textbf{永不缺货}：每个数量级都有新鲜素数入库（欧几里得的保证兑现为数据）；
\item \textbf{密度一路走低}：从 $40\%$ 掉到 $2.47\%$，而且看起来还要继续掉——素数越来越"稀"；
\item \textbf{掉得有节奏}：数量级每升一级（$x$ 乘 $10$），占比大致缩到原来的 $\tfrac45$ 到 $\tfrac9{10}$ 之间，似乎奔某个确定的比例去。
\end{enumerate}
第 $2$、$3$ 条之间的张力就是本章的悬念：\textbf{"越来越稀"与"按什么速度稀"是两个问题}。前者已见，后者要找一条\textbf{标尺}。

\section{找标尺：为什么偏偏是对数}

猜标尺是个技术活。第一个候选：占比缩到原来的固定比例？$25\% \to 16.8\% \to 12.3\%$，比例是 $0.672, 0.732, \dots$——不固定，否决。第二个候选来自第 19 章的直觉：占比的倒数 $\dfrac{x}{\pi(x)}$（"平均多少个数里出一个素数"）：

\begin{center}
\begin{tabular}{ccc}
\toprule
$x$ & $x/\pi(x)$（平均间距） & $\ln x$ \\
\midrule
$100$ & $4.0$ & $4.605$ \\
$1\,000$ & $5.95$ & $6.908$ \\
$10^4$ & $8.14$ & $9.210$ \\
$10^6$ & $12.74$ & $13.816$ \\
$10^9$ & $19.67$ & $20.723$ \\
$10^{18}$ & $40.42$ & $41.447$ \\
\bottomrule
\end{tabular}
\end{center}

两列数字\textbf{并肩前行}，始终略低于 $\ln x$，且越靠后贴得越近（$10^{18}$ 处只差 $2.4\%$）。换算成一条猜想：

\begin{caixiang}[素数定理]
$x$ 很大时，
\[
\pi(x) \;\approx\; \frac{x}{\ln x},
\qquad \text{即} \qquad
\lim_{x\to\infty} \frac{\pi(x)}{x/\ln x} = 1 .
\]
\end{caixiang}

\textbf{为什么标尺偏偏是 $\ln$？}给一个"筛子直觉"（不是证明，是让人睡得着觉的理由）：在 $x$ 附近随机抓一个整数，它是素数的概率有多大？它得先逃过 $2$ 的筛（幸存率 $\tfrac12$，因为一半的数被 $2$ 整除）、再逃过 $3$ 的筛（幸存率再乘 $\tfrac23$）、逃过 $5$（乘 $\tfrac45$）……对每个不超过 $\sqrt x$ 的素数 $p$，都要缴纳 $\left(1 - \tfrac1p\right)$ 的"过路费"。总幸存率：
\[
\prod_{p \le \sqrt x}\left(1 - \frac1p\right)
\;\;\overset{?}{\sim}\;\; \frac{1}{\ln x} .
\]
（"$\sim$"是"渐近相等"的记号。）这个乘积到底等于什么？第 22 章将证明它恰好是 $\tfrac{1}{\ln x}$ 量级——欧拉的齿轮会咬合出对数。\textbf{素数密度 $= \tfrac1{\ln x}$，累计起来 $\pi(x) \approx \int \tfrac{dt}{\ln t} \approx \tfrac{x}{\ln x}$}：对数就这样从"连乘的筛子"里长了出来。第 23 章把这条直觉打磨成定理。

\begin{cidian}
"$\approx$"的正式含义（土名"越晚越贴"$\to$ \textbf{渐近}）：两者之比的落脚点是 $1$。它\textbf{不}保证逐点误差小——$x/\ln x$ 在 $x = 10^9$ 处仍低估约 $5.1\%$，但比值死心塌地奔 $1$。粗估 $x/\ln x$ 与精估 $\mathrm{li}(x)$ 的区别，第 23 章专讲。
\end{cidian}

\section{高斯的传说}

1792 年，十五岁的高斯得到一本对数表，表尾附了一张素数表。别的少年用来查题，他干了件体力活：把每千个整数一段的素数数了一遍又一遍，然后宣布——"每千个数里的素数个数，约等于 $\tfrac{1000}{\ln x}$"（$x$ 是该段的位置）。顺手把素数定理猜了出来。

这个故事的恐怖之处有三层。第一，\textbf{他猜的是"密度 $\tfrac1{\ln x}$"而不是"$\pi(x)\approx x/\ln x$"}——前者是微分式的眼光（局部密度），后者是积分式的结论，十五岁的直觉直接切进了更深的层次（密度累计即 $\mathrm{li}$，第 23 章）。第二，猜完他不撒手，\textbf{亲手把素数数到三百万}验证（无计算机时代，数月之功）。第三，他一生没发表这条猜想——手稿里躺着，直到 1849 年给恩克的信里才提起（"我记不清是否早就告诉过你"）。数学史上最富的人随手埋了张彩票。

\begin{lishi}
素数定理的证明史是一台跨世纪接力机。勒让德 1798 年猜出 $\pi(x) \approx \tfrac{x}{\ln x - 1.08366}$（经验公式，常数纯靠拟合）；切比雪夫 1852 年用初等方法证明 $\pi(x)$ 被夹在 $0.92\,\tfrac{x}{\ln x}$ 与 $1.11\,\tfrac{x}{\ln x}$ 之间——离终点一步之遥却停在那里（初等工具的极限）；1896 年，阿达马与德拉瓦莱-普桑\textbf{同一年、各自独立}、用同一件刚出厂的重型武器（第 24 章的 $\zeta$ 函数在复平面的性质）冲线。一条关于"数整数"的定理，要动用"连续曲线"的复分析：这是本书最震撼的跨界事件，也是"解析数论"四个字的含义——\textbf{用分析（微积分与复变）的显微镜，看整数世界的细胞}。
\end{lishi}

\begin{dongshou}
\begin{itemize}
  \yisi 补普查表的行：$x = 10^5$，$\pi(x) = 9592$。计算 $\ln 10^5 = 5\ln10 \approx 11.513$，$x/\ln x \approx 8686$，比值 $\pi\cdot\ln x/x \approx 1.104$。把这一行插进表里，确认比值序列 $1.15, 1.16, 1.13, \dots$ 单调下滑奔 $1$。
  \yier 用 $\ln10 \approx 2.3026$ 推算 $\ln 10^8$，再用 $x/\ln x$ 预估一亿以内的素数个数（真值 $5{,}761{,}455$；你的预估 $\approx 5.43\times10^6$，误差 $5.7\%$——然后算"每几个数里一个素数"：$\tfrac{10^8}{5.76\times10^6} \approx 17.4$，与 $\ln 10^8 = 18.4$ 对照）。
  \yier 在 $1$ 到 $200$ 里数素数（$46$ 个），分前后两半统计：$1$--$100$ 有 $25$ 个（密度 $25\%$），$101$--$200$ 有 $21$ 个（$21\%$）——下降肉眼可见。再用密度 $\tfrac1{\ln x}$ 预测 $201$--$300$：取段中 $x = 250$，$\tfrac{100}{\ln 250} \approx \tfrac{100}{5.52} \approx 18$ 个；真值 $16$ 个，误差约两个。全局累计 $\pi(x)$ 的预测越来越准（正文表格），而\textbf{小区间}的预测永远"颠簸"几个——素数的局部波动是顽疾，这份颠簸正是第 23 章 $\mathrm{li}(x)$ 与第 24 章零点要治理的对象。
\end{itemize}
\end{dongshou}

\begin{shaonao}
孪生素数（第 4 章）也有"人口普查"：$x$ 以内孪生素数对数约为 $1.32\,\dfrac{x}{(\ln x)^2}$（哈代--李特尔伍德猜想，1923；注意分母是\textbf{对数的平方}——孪生素数比单个素数稀少一个"对数量级"）。用 $x = 10^6$ 检验：$1.32\times10^6/(13.816)^2 \approx 6920$ 对；真值 $8169$ 对，比值 $1.18$。再看 $x = 10^9$：预测 $\approx 3.09\times10^6$ 对，真值 $3.424\times10^6$，比值 $1.11$——同样在缓慢奔 $1$。做两件事：(1) 把两个比值与正文表格里 $\pi(x)$ 的比值（同位置 $1.08, 1.05$）并排放，观察"猜想越难，比值收敛越慢"；(2) 解释为什么分母是 $(\ln x)^2$：孪生素数像"素数密度过两次筛"的联合事件，$\tfrac1{\ln x}$ 自乘一次。这道题让你亲手用初等算术摸到了未解之谜的数值面——\textbf{猜想的"猜"，是有结构的猜}。
\end{shaonao}

\begin{chengshi}
本章完全没碰证明——素数定理的证明需要 $\zeta$ 函数与复分析，超过全书工具，这是全书最大的一处"地毯"，地毯下的第一块砖在第 22 章（欧拉乘积）铺。"筛子直觉"里 $\prod(1-\tfrac1p) \sim \tfrac1{\ln x}$（梅尔滕斯第三定理，1874）同样是专业成果。$10^{18}$ 行的 $\pi$ 值来自 2015 年的分布式计算（Büthe 等人验证至 $10^{19}$），这类计算本身就是素数研究的日常。
\end{chengshi}

预报员 $x/\ln x$ 为什么越晚越准、它和更准的 $\mathrm{li}(x)$ 差多少？在揭开谜底之前，需要先备料：研究一种\textbf{加不拢的无穷}——它是素数世界的温度计。
